\begin{proposition}[Proved extension of the marked torus period]\hypertarget{tag-01GD}{}\label{cdp:st:D.2.6}
\StatementRecord{cdp:st:D.2.6}{proved}{01GD}{\SeriesRef{cdp:st:D.2.3}, \SeriesRef{cdp:eq:D.20}}


Let $h:U\to\Delta$ be a proper holomorphic map from a smooth
complex threefold, smooth with connected two-torus fibres over
$\Delta^*$. Suppose that its integral $H^1$ monodromy is the
identity. The marked period map
\[
 \Delta^*\longrightarrow
 \mathcal T_2
 :=\{W\in\operatorname{Gr}(2,H^1(U_t;\C)):
                   W\oplus\overline W=H^1(U_t;\C)\}
\]
extends holomorphically across zero with value in $\mathcal T_2$.
In particular the relative Albanese group family over $\Delta^*$
extends to a smooth proper group family of complex two-tori over
$\Delta$. This assertion does not assert an extension of the
Albanese map from the given proper model $U$.
\end{proposition}


\begin{proof}

\begin{enumerate}[label={$\langle1\rangle$\arabic*.},leftmargin=*]
\item An embedded resolution of the central fibre changes neither
the smooth fibres nor their marking, so we may use the relative
logarithmic complex of \ref{cdp:st:D.2.3} with arbitrary multiplicities.
Identity on $H^1$ gives identity on $H^2=\bigwedge^2H^1$.
Consequently its logarithmic Gauss--Manin bundle has zero
residue and a holomorphic flat frame. Let $\Omega$ be a
generator of $h_*\mathscr K^2$, as in \ref{cdp:st:D.2.3}.

\item The second case of that proof, including the exact maps
(D.18)--(D.19), did not use any geometric-genus assumption.
It shows that $\Omega$ cannot vanish on the whole reduced
central support: otherwise its divided cohomology class would
have bounded periods while its positive $L^2$ volume diverges.
Thus $\Omega|_T$ is nonzero on some normalized component $T$.
The first case of that proof shows, without a geometric-genus
assumption, that every conductor residue of every component
restriction is zero, because a nonzero residue would give the
same divergent positive volume. The nonzero restriction
$\Omega|_T$ is therefore a holomorphic two-form on $T$.

\item Choose a point of $T$ outside the conductor and the zeros of
this two-form. In a small chart the map is $t=z^m$ and
\[
 \Omega=a(z,\xi,\eta)\,d\xi\wedge d\eta,
 \qquad a(0,0,0)\ne0.
\]
Integrating on one sheet over a fixed small
$(\xi,\eta)$-polydisc yields a constant $c>0$ such that
\[
 \int_{U_t}\Omega_t\wedge\overline{\Omega_t}\ge c
 \qquad(0<|t|\ll1).
\tag{D.20}\hypertarget{tag-0197}{}\label{cdp:eq:D.20}
\]
The flat cohomology class $v(t)=[\Omega_t]$ is holomorphic
across zero. Its cup pairing with its conjugate equals the
integral in \eqref{cdp:eq:D.20}, hence
\[
 v(0)\wedge\overline{v(0)}>0
 \quad\hbox{in }H^4(U_t;\R)\simeq\R,
\]
using the complex orientation. In particular $v(0)\ne0$.

\item Write $V=H^1(U_t;\C)$ with its fixed integral marking.
For $t\ne0$, $v(t)$ is a decomposable element of
$\bigwedge^2V$, spanning $\bigwedge^2H^{1,0}(U_t)$.
The decomposable cone is the closed algebraic set defined by
the Pl\"ucker equations, so $v(0)$ is decomposable. For a
nonzero decomposable $v$, the exact inverse to the Pl\"ucker
map is
\[
 W(v)=\{\alpha\in V:\alpha\wedge v=0\}.
\]
This kernel has dimension two and varies holomorphically near
$v(0)$. The nonzero pairing
$v(0)\wedge\overline{v(0)}$ says exactly that
$W(v(0))\cap\overline{W(v(0))}=0$. The extended period
therefore lies in $\mathcal T_2$, as claimed.

\item For completeness the extending group family is explicit.
Let $\mathcal W\subset V\otimes\OO_\Delta$ be the extended
rank-two subbundle $W(v(t))$, and let
$\Lambda=H_1(U_t;\Z)$ with its fixed marking. Evaluation of
cohomology on integral cycles gives a holomorphic embedding
of the constant lattice $\Lambda$ in $\mathcal W^*$.
The direct-sum relation with the conjugate makes this a full
real lattice in every fibre, uniformly on compact subsets of
the disc. Hence
\[
 \mathcal A=\mathcal W^*/\Lambda\longrightarrow\Delta
\]
is a smooth proper family of complex Lie groups with
two-torus fibres. On $\Delta^*$, integration of holomorphic
one-forms identifies each group fibre with the Albanese torus
of $U_t$, and these identifications vary holomorphically.
This is the asserted exact isomorphism of relative group
families. Extending a map from the original model $U$ to
$\mathcal A$ requires a further argument; it does not follow
merely from this period construction.

\end{enumerate}
\end{proof}